\chapter{Описание статического анализа} \label{ch3}

Данная глава посвящена выбору статических анализаторов, их настройке, результатам
проверки кода приложения и сопоставлению найденных предупреждений с CWE и БДУ
ФСТЭК России.

\section{Выбор анализаторов и источники основной информации} \label{ch3:sec1}

Для дальнейшей работы выбраны два анализатора:
\begin{enumerate}
    \item SonarJS-совместимые правила ESLint как локальный вариант правил SonarSource;
    \item ESLint Security как специализированный анализатор потенциальных проблем безопасности JavaScript-кода.
\end{enumerate}

SonarQube выбран как анализатор из списка ``A''. В официальной документации
указано, что анализ JavaScript, TypeScript и CSS доступен в SonarQube Server и
Community Build [5]. В локальном проекте использованы
SonarJS-совместимые правила через \texttt{eslint-plugin-sonarjs}: репозиторий
плагина указывает, что новые версии делают правила SonarJS доступными
пользователям ESLint [4]. ESLint Security выбран как второй
анализатор, так как он содержит правила для поиска небезопасных регулярных
выражений, object injection и опасного использования файловой системы
[3].

Ссылки на источники основной информации по анализаторам приведены в
\taref{tab:analyzer-links}.

\noindent
\begingroup
\centering
\small
\begin{longtable}[c]{|p{4cm}|p{5.5cm}|p{5.5cm}|}
    \caption{Источники информации по анализаторам}
    \label{tab:analyzer-links}
    \\
    \hline
    Анализатор & Основная информация & Страница с правилами / уязвимостями \\ \hline
    \endfirsthead
    \hline
    Анализатор & Основная информация & Страница с правилами / уязвимостями \\ \hline
    \endhead
    \hline
    SonarQube / SonarJS & Документация SonarSource по анализу JavaScript [5] & Репозиторий SonarJS-правил для ESLint [4] \\ \hline
    ESLint Security & Страница пакета \texttt{eslint-plugin-security} [3] & Список правил пакета \texttt{eslint-plugin-security} [3] \\ \hline
\end{longtable}
\normalsize
\endgroup

\section{Основные вехи установки и настройки} \label{ch3:sec2}

Установка выполнялась через npm:
\begin{lstlisting}[breaklines=true]
npm install eslint @eslint/js eslint-plugin-sonarjs eslint-plugin-security --save-dev
\end{lstlisting}

Настройка выполнена в файле \texttt{eslint.config.js}. В конфигурацию включены
правила:
\begin{itemize}
    \item \texttt{security/detect-unsafe-regex};
    \item \texttt{security/detect-object-injection};
    \item \texttt{security/detect-non-literal-fs-filename};
    \item \texttt{sonarjs/slow-regex};
    \item \texttt{sonarjs/no-duplicated-branches};
    \item \texttt{sonarjs/cognitive-complexity}.
\end{itemize}

Файл конфигурации задаёт область анализа: каталоги \texttt{src}, \texttt{fuzz} и
\texttt{scripts}. В список глобальных идентификаторов добавлены \texttt{console},
\texttt{module}, \texttt{require}, \texttt{process}, \texttt{Buffer} и
\texttt{\_\_dirname}, так как проект написан в формате CommonJS и запускается в
Node.js.

Запуск анализа:
\begin{lstlisting}[breaklines=true]
npm run lint:security
npm run static:summary
\end{lstlisting}

Первая команда формирует JSON-отчёт \texttt{reports/static/eslint-security.json}.
Вторая команда преобразует его в CSV-сводку \texttt{reports/static/static-summary.csv}.

При настройке возникла дополнительная зависимость \texttt{ts-api-utils}, требуемая
пакетом \texttt{eslint-plugin-sonarjs}. Она была установлена через npm. После этого
оба анализатора запускались одной командой ESLint и формировали общий JSON-отчёт.

\section{Основные виды статического анализа} \label{ch3:sec3}

Выбранные анализаторы выполняют следующие виды проверок:
\begin{itemize}
    \item поиск потенциально небезопасных регулярных выражений;
    \item поиск возможной object injection при обращении к объекту или массиву по динамическому ключу;
    \item поиск операций файловой системы с нелитеральным именем файла;
    \item поиск неиспользуемых переменных;
    \item поиск сложных или дублирующихся ветвей;
    \item поиск проблем сопровождаемости кода.
\end{itemize}

Для классификации дефектов использованы Common Weakness Enumeration MITRE [7]
и Банк данных угроз безопасности информации ФСТЭК России [6].

Классы проверок по каждому анализатору приведены в \taref{tab:analysis-types}.

\noindent
\begingroup
\centering
\small
\begin{longtable}[c]{|p{4cm}|p{10.5cm}|}
    \caption{Виды статического анализа}
    \label{tab:analysis-types}
    \\
    \hline
    Анализатор & Виды анализа \\ \hline
    \endfirsthead
    \hline
    Анализатор & Виды анализа \\ \hline
    \endhead
    \hline
    ESLint Security & Поиск небезопасных регулярных выражений, динамических путей файлов, потенциальной object injection, небезопасных вызовов API \\ \hline
    SonarJS & Поиск дефектов сопровождаемости, дублирующихся ветвей, сложных условий, медленных регулярных выражений, проблем читаемости \\ \hline
\end{longtable}
\normalsize
\endgroup

\section{Итоговая таблица результатов статического анализа} \label{ch3:sec4}

Итоговая таблица по предупреждениям, релевантным внесённым ошибкам и анализируемым
рискам, приведена в \taref{tab:static-results}.

\noindent
\begingroup
\centering
\small
\begin{longtable}[c]{|p{3.6cm}|p{2.8cm}|p{2.8cm}|p{2.2cm}|p{3.1cm}|}
    \caption{Комплексная информация по предупреждениям статических анализаторов}
    \label{tab:static-results}
    \\
    \hline
    Файл & Анализатор / правило & CWE & БДУ ФСТЭК & Описание предупреждения \\ \hline
    \endfirsthead
    \hline
    Файл & Анализатор / правило & CWE & БДУ ФСТЭК & Описание предупреждения \\ \hline
    \endhead
    \hline
    \texttt{datasetNameParser.js:2} & \texttt{security/detect-unsafe-regex} & CWE-1333 & УБИ.038 & Небезопасное регулярное выражение может привести к отказу в обслуживании \\ \hline
    \texttt{datasetNameParser.js:2} & \texttt{sonarjs/slow-regex} & CWE-1333 & УБИ.038 & Регулярное выражение уязвимо к сверхлинейному времени из-за backtracking \\ \hline
    \texttt{jsonReader.js:44} & \texttt{security/detect-non-literal-fs-filename} & CWE-22 & УБИ.028 & Использование пути к файлу, полученного не из литерала \\ \hline
    \texttt{resultWriter.js:20} & \texttt{security/detect-non-literal-fs-filename} & CWE-22 & УБИ.028 & Запись файла по динамическому пути \\ \hline
    \texttt{manacherOdd.js:22} & \texttt{security/detect-object-injection} & CWE-94 & УБИ.028 & Динамическое обращение по индексу; для массива алгоритма является допустимым, но анализатор помечает как риск \\ \hline
    \texttt{jsonReader.js:37} & \texttt{no-unused-vars} & CWE-563 & УБИ.025 & Переменная \texttt{error} объявлена, но не используется \\ \hline
\end{longtable}
\normalsize
\endgroup

Ошибка CWE-1333 относится к классу слабостей, связанных с неэффективной
сложностью регулярного выражения [8]. Ошибки CWE-22 и CWE-94 относятся к
распространённым классам дефектов безопасности и требуют проверки источника
входных данных. В данном приложении часть предупреждений
\texttt{detect-object-injection} является ложноположительной, так как массивы
радиусов индексируются числовыми индексами алгоритма, а не пользовательскими
ключами.

Отдельно следует отметить предупреждение \texttt{sonarjs/no-duplicated-branches}
для модуля \texttt{percentInputDecoder.js}. Оно относится не к аварийному дефекту,
найденному фаззером, а к качеству кода: две ветви \texttt{case 0} и \texttt{case 1}
возвращают одинаковое значение \texttt{"short"}. Такое предупреждение не является
критическим, но показывает, что статический анализатор выявляет не только
уязвимости, но и проблемы сопровождаемости.

\section{Коды ошибок и сопоставление с CWE и БДУ ФСТЭК} \label{ch3:sec5}

Сопоставление найденных предупреждений с классификаторами приведено в
\taref{tab:cwe-bdu}. Для БДУ ФСТЭК использованы близкие по смыслу классы угроз,
так как БДУ описывает угрозы безопасности информации, а CWE описывает слабости
программного обеспечения.

\noindent
\begingroup
\centering
\small
\begin{longtable}[c]{|p{3.2cm}|p{3.6cm}|p{3.2cm}|p{4.2cm}|}
    \caption{Сопоставление предупреждений с CWE и БДУ ФСТЭК}
    \label{tab:cwe-bdu}
    \\
    \hline
    Правило & CWE & БДУ ФСТЭК & Обоснование сопоставления \\ \hline
    \endfirsthead
    \hline
    Правило & CWE & БДУ ФСТЭК & Обоснование сопоставления \\ \hline
    \endhead
    \hline
    \texttt{detect-unsafe-regex} & CWE-1333 & УБИ.038 & Риск отказа в обслуживании из-за сверхлинейного времени регулярного выражения \\ \hline
    \texttt{slow-regex} & CWE-1333 & УБИ.038 & Аналогичный риск ReDoS через backtracking \\ \hline
    \texttt{detect-non-literal-fs-filename} & CWE-22 & УБИ.028 & Риск обращения к файлу по неконтролируемому пути \\ \hline
    \texttt{detect-object-injection} & CWE-94 & УБИ.028 & Потенциальное использование пользовательского ключа для обращения к объекту \\ \hline
    \texttt{no-unused-vars} & CWE-563 & УБИ.025 & Лишняя переменная может скрывать незавершённую обработку исключения \\ \hline
\end{longtable}
\normalsize
\endgroup

CWE-1333 связан с угрозами отказа в обслуживании и учитывается в материалах MITRE
по опасным ошибкам программного обеспечения [8, 9]. Для остальных предупреждений
прямая принадлежность к Top-25 зависит от года рейтинга и контекста применения;
в рамках данной работы они рассматриваются как значимые, но требующие ручной
проверки.

\section{Обсуждение результатов и сравнение анализаторов} \label{ch3:sec6}

ESLint Security оказался более полезным для поиска дефектов безопасности в
JavaScript-коде. Он обнаружил небезопасное регулярное выражение, динамические пути
файлов и потенциальную object injection. Его преимущество --- простая установка и
понятные rule id. Недостаток --- значительное число предупреждений требует ручной
проверки контекста.

SonarJS-совместимые правила дополнили анализ предупреждением \texttt{slow-regex}.
Они полезны для проверки качества и сопровождаемости кода, а также для выявления
части дефектов производительности. В контексте данного приложения SonarJS точнее
описал риск регулярного выражения как сверхлинейное время из-за backtracking.

Сравнение анализаторов приведено в \taref{tab:analyzer-comparison}.

\noindent
\begingroup
\centering
\small
\begin{longtable}[c]{|p{4cm}|p{5.5cm}|p{5.5cm}|}
    \caption{Сравнение удобства и эффективности анализаторов}
    \label{tab:analyzer-comparison}
    \\
    \hline
    Критерий & ESLint Security & SonarJS-правила \\ \hline
    \endfirsthead
    \hline
    Критерий & ESLint Security & SonarJS-правила \\ \hline
    \endhead
    \hline
    Установка & Через npm, без сервера & Через npm, без сервера \\ \hline
    Типы предупреждений & Безопасность JavaScript: regex, object injection, fs path & Качество, сопровождаемость, regex-производительность \\ \hline
    Полезность для проекта & Высокая & Средняя \\ \hline
    Ложноположительные срабатывания & Есть на индексах массивов & Меньше на выбранных правилах \\ \hline
    Удобство отчёта & JSON-отчёт и CSV-сводка & Включается в тот же JSON-отчёт ESLint \\ \hline
\end{longtable}
\normalsize
\endgroup

В контексте поиска ошибок в данном коде наиболее эффективной оказалась связка
инструментов: ESLint Security обнаруживает больше потенциальных рисков безопасности,
а SonarJS уточняет качество отдельных конструкций. При этом ни один статический
анализатор не нашёл напрямую отсутствие \texttt{try/catch} вокруг
\texttt{decodeURIComponent}; этот дефект был обнаружен именно фаззингом. Это
подтверждает необходимость совместного применения статического и динамического
анализа.

\section{Выводы} \label{ch3:conclusion}

Статический анализ выявил предупреждения, связанные с небезопасным регулярным
выражением, динамическими путями файлов, потенциальной object injection и
неиспользуемой переменной. Наиболее значимым предупреждением для внесённых ошибок
является \texttt{security/detect-unsafe-regex} и \texttt{sonarjs/slow-regex}, так
как оно указывает на риск отказа в обслуживании при обработке пользовательской
строки. Результаты статического анализа дополняют фаззинг: фаззер обнаружил
необработанные исключения на конкретных входах, а статические анализаторы указали
на потенциальные классы дефектов без выполнения программы.
